\documentclass[conference]{IEEEtran}
\usepackage{cite}
\usepackage{booktabs}
\title{Cross-Script Offensive Language Detection for Roman Urdu: Intermediate Report}
\author{\IEEEauthorblockN{Group 18}\IEEEauthorblockA{CS 5316: NLP Theory and Applications}}
\begin{document}\maketitle
\begin{abstract}
Roman Urdu offensive-language detection must handle informal transliteration, spelling variation, and code-switching. We implement a character n-gram TF--IDF and linear SVM baseline on the coarse-grained Roman Urdu Hate Speech and Offensive Language Dataset (RUHSOLD). On the supplied test split, it reaches 89.02\% accuracy and 0.8894 macro-F1. Our proposed extension is Cross-Script Alignment Prompting (CSAP), which provides validated Roman Urdu--Urdu correspondences to a prompted language model. Baseline evaluation is complete; CSAP inference and evaluation remain planned work, and no novelty improvement is claimed.
\end{abstract}
\begin{IEEEkeywords}Roman Urdu, offensive language detection, text classification, transliteration, prompting\end{IEEEkeywords}
\section{Introduction}
Roman Urdu social-media text is written in informal Latin transliteration, often with variable spellings and English code-switching. This creates a normalization challenge for offensive-language detection. We implement a reproducible supervised baseline and propose CSAP, which adds a small, curated cross-script lexicon to an LLM classification prompt. The project is informed by prompt-based Urdu detection research~\cite{maab2026}; our completed experiment uses the locally available RUHSOLD release~\cite{rizwan2020}, so it is a project baseline rather than a complete reproduction of the anchor study.
\section{Methodology}
\subsection{Data and baseline}
We use RUHSOLD Task 1, a binary coarse-grained classification set with supplied train, validation, and test TSVs. The dataset definitions map label 0 to abusive/offensive and label 1 to normal. Split sizes are 7,209, 801, and 2,003. We retain the official split and raw text, removing only malformed or empty rows. The baseline is character TF--IDF (n-grams 2--5) with LinearSVC ($C=1$), implemented in scikit-learn. The validation split is reserved and was not used for tuning. Task 2 is excluded because its local validation TSV duplicates its training TSV.
\subsection{Planned novelty and evaluation}
CSAP will be tested against zero-shot and fixed four-example few-shot prompts, with alignment conditions $N\in\{5,15\}$; $N=0$ is the few-shot control. Before inference, we will manually verify each mapping against Urdu spelling and usage, freeze the lexicon and demonstrations, and publish them with the run configuration. We will select a deterministic, label-stratified subset from the held-out test split, run every prompt condition on the same examples with one model and decoding configuration, and save each response immediately for resumability. We will report macro-F1, per-class scores, accuracy, coverage/failures, latency, and prompt length. Results will be interpreted as a pilot if API limits constrain sample size. No LLM inference has yet been completed. The progress-checkpoint instructions allow either an initial novelty implementation or a clear description of the planned work. We follow the latter: to avoid hosted API quota constraints, we will run the prompt conditions locally through Ollama, beginning with a quantized Qwen 3.5 4B model. This is a concrete plan, not an experimental finding.
\section{Results}
The character TF--IDF + LinearSVC baseline achieves 89.02\% accuracy, 0.8894 macro-F1, and 0.8900 weighted F1 on the 2,003-example test split. For abusive/offensive, precision, recall, and F1 are 0.8948, 0.8660, and 0.8802; for normal they are 0.8864, 0.9112, and 0.8986. The confusion matrix, with gold and predicted order [abusive/offensive, normal], is $[[808,125],[95,975]]$. No CSAP or LLM results are available yet. Metrics and per-example predictions are saved in \texttt{results/}; the runnable baseline notebook is in \texttt{notebooks/}.
\section{Next steps}
The supervised baseline and held-out evaluation are complete. Remaining work is to validate and freeze the alignment lexicon and demonstrations, implement resumable LLM logging, run matched prompt conditions locally through Ollama (initial model: quantized Qwen 3.5 4B), and analyze errors and spelling-variation robustness. The current evidence establishes a classical reference point only; whether CSAP improves detection remains an open question.
\begin{thebibliography}{2}
\bibitem{maab2026} I. Maab, U. Haider, and J. Yamagishi, ``Prompt-driven Detection of Offensive Urdu Language using Large Language Models,'' in \emph{Proc. EACL}, 2026, pp. 4302--4327.
\bibitem{rizwan2020} H. Rizwan, M. H. Shakeel, and A. Karim, ``Hate-speech and offensive language detection in Roman Urdu,'' in \emph{Proc. EMNLP}, 2020, pp. 2512--2522.
\end{thebibliography}
\end{document}
